\section{保号、有限首段与零点定位}\label{cd:limits-sign}
\begin{examtip}
极限给出的是最终成立的固定精度估计。利用这种估计，可以把数列的最值问题压到有限首段，也可以在开区间或无界区间上选出适合介值定理的两个有限端点。
\end{examtip}

\subsection{极限附近的固定间隔与有限首段}
\begin{method}{先固定间隔，再截取尾部}
\textbf{适用条件} 数列$a_n\to L$，且已找到$a_p>L$或$a_q<L$。\par
\textbf{步骤} 用已知项与$L$之间的距离选定正数$\varepsilon$；将所有充分靠后的项收进$(L-\varepsilon,L+\varepsilon)$；最后只比较有限首段。\par
\textbf{疑难点} $\varepsilon$必须先固定。不能把“对任意固定$\varepsilon$最终成立”直接替换成$a_n-L<\frac{1}{n}$。
\end{method}
\begin{theorem}{收敛数列的最值取得判据}\label{cd:finite-extrema}
设$a_n\to L$。若存在$p$使$a_p>L$，则数列取得最大值；若存在$q$使$a_q<L$，则取得最小值。特别地，若数列在$L$两侧各有至少一项，则最大值、最小值均能取得。
\end{theorem}
\begin{proof}
令$\varepsilon=\frac{a_p-L}{2}>0$。存在$N\ge p$，使$n>N$时
\[
 a_n<L+\varepsilon=\dfrac{a_p+L}{2}<a_p.
\]
因此有限集合$\{a_1,\ldots,a_N\}$的最大值就是全数列最大值。对最小值取$\varepsilon=\frac{L-a_q}{2}$，同理可证。
\end{proof}
\begin{example}{考研风格｜有最值与仅有上下确界}{难度：$\star\star$}\label{cd:ex:finite-extrema}
\textbf{题目} 判断$a_n=\frac{\sin n}{n}$是否取得最大值和最小值；再与$b_n=\frac{1}{n}$比较。\par
\textbf{思路} 无须先找到最大项和最小项，只需找到极限两侧的项。\par
\textbf{详细解答} 由$|a_n|\le\frac{1}{n}$得$a_n\to0$。又$\sin1>0$，而$\pi<4<\frac{3\pi}{2}$，故$a_1>0$、$a_4<0$。应用定理\ref{cd:finite-extrema}，两种最值都能取得。直接证明时可取
\[
 \varepsilon=\dfrac12\min\{a_1,-a_4\}>0.
\]
当$n>N\ge4$时，$-\varepsilon<a_n<\varepsilon$，尾项严格小于$a_1$且大于$a_4$，故所有最值都可在前$N$项中找。对$b_n$，最大值为$b_1=1$；每项为正且趋0，故下确界为0，但没有最小值。\par
\textbf{方法点评} “收敛所以有界”只能保证上下确界存在；极限两侧的具体项才把确界提升为实际取得的最值。\par
\textbf{疑难点与变式} 结论只证明最值存在，没有声称$a_1$或$a_4$一定就是极值项。若只在极限一侧有项，另一种最值仍须另查。
\end{example}

\subsection{平移差与相等函数值}
\begin{method}{将相等函数值改为连续差的零点}
\textbf{适用条件} 要证明$f(t)=f(t+h)$，且两端点条件具有平移对称。\par
\textbf{步骤} 构造$G(t)=f(t)-f(t+h)$，在其定义区间端点计算符号；端点乘积非正时用介值定理。\par
\textbf{疑难点} 非正乘积允许零点在端点；只有严格异号才能保证零点在开区间。
\end{method}
\begin{example}{考研风格｜相隔定长的两个等高点}{难度：$\star\star$}\label{cd:ex:translate}
\textbf{题目} 设$a>0$，$f\in C[0,2a]$且$f(0)=f(2a)$，证明存在$t\in[0,a]$，使$f(t)=f(t+a)$。说明$t$能否总取在$(0,a)$。\par
\textbf{思路} 比较从左半区间和右半区间得到的同一差。\par
\textbf{详细解答} 令$G(t)=f(t)-f(t+a)$，则$G\in C[0,a]$，且
\[
 G(a)=f(a)-f(2a)=f(a)-f(0)=-G(0).
\]
若$G(0)=0$，取$t=0$；否则$G(0)G(a)=-G(0)^2<0$，介值定理给$t\in(0,a)$满足$G(t)=0$。开区间结论不总成立：取$f(x)=\sin(\frac{\pi x}{a})$，端点$f(0)=f(2a)=0$，但对$0<t<a$，
\[
 f(t)-f(t+a)=2\sin(\dfrac{\pi t}{a})>0,
\]
所以等高点只对应$t=0,a$。\par
\textbf{方法点评} 用差函数替代同时寻找两个点，可将几何等高问题变成一个未知数的零点问题。\par
\textbf{疑难点与变式} 若额外有$f(0)\ne f(a)$，端点严格异号，便能保证$t\in(0,a)$。
\end{example}

\subsection{归一化符号与无界区间零点}
\begin{method}{先在两端归一化，再截成闭区间}
\textbf{适用条件} $H$在$(0,\infty)$连续，能够分别在$0^+$和$+\infty$找到其符号。\par
\textbf{步骤} 用正的尺度因子消去趋零或增长的主量；由非零极限选出有限$x_1<x_2$使$H(x_1)H(x_2)<0$；最后在$[x_1,x_2]$使用介值定理。\par
\textbf{疑难点} 无界区间的两端不是可直接代入的函数值，必须实际选出有限端点。
\end{method}
\begin{example}{考研风格｜局部三阶信息强制出现零点}{难度：$\star\star\star$}\label{cd:ex:unbounded-root}
\textbf{题目} 设$f$在$(0,\infty)$连续有界，且
\[
 \lim_{x\to0^+}\dfrac{f(x)}{\sin x-x}=12.
\]
证明方程$f(x)+x^3=0$在$(0,\infty)$至少有一个实根。\par
\textbf{思路} 在0附近比较三阶主项，在无穷远用有界性压低$f$的阶。\par
\textbf{详细解答} 记$H(x)=f(x)+x^3$。由于$\frac{\sin x-x}{x^3}\to-\frac{1}{6}$，
\[
 \lim_{x\to0^+}\dfrac{H(x)}{x^3}
 =12\left(-\dfrac16\right)+1=-1.
\]
故存在$\delta>0$，使$0<x<\delta$时$H(x)<0$。另一方面，由$|f(x)|\le M$，
\[
 \left|\dfrac{f(x)}{x^3}\right|\le\dfrac{M}{x^3}\to0,
 \qquad \dfrac{H(x)}{x^3}\to1\quad(x\to+\infty).
\]
选$x_1\in(0,\delta)$，再取足够大的$x_2>x_1$使$H(x_2)>0$。$H$在$[x_1,x_2]$连续，介值定理给出所需零点。\par
\textbf{方法点评} 极限只负责找到符号，连续性负责连接这两个符号；两者的职责应分别写清。\par
\textbf{疑难点与变式} 本题没有单调性或导数信息，不能进一步断言根唯一。将12换成参数$c$时，上述论证在$c>6$仍成立。
\end{example}
